% !TEX root = ../DoAn.tex
\documentclass[../DoAn.tex]{subfiles}
\begin{document}

Phụ lục này cung cấp các lệnh vận hành chính dùng để biên dịch, khởi động, kiểm tra và dừng
hệ thống trong môi trường local và môi trường AWS EC2 dạng production-like.

\section*{A.1. Yêu cầu hệ thống}

\begin{itemize}
    \item \textbf{Java:} OpenJDK 21.
    \item \textbf{Docker:} Docker Desktop 27+ hoặc Docker Engine kèm Docker Compose plugin.
    \item \textbf{Maven:} Không cần cài riêng vì dự án dùng Maven Wrapper \texttt{./mvnw}.
    \item \textbf{Node.js:} Node.js 20+ cho frontend Next.js.
    \item \textbf{RAM:} Tối thiểu 12GB khả dụng cho Docker, khuyến nghị 16GB.
    \item \textbf{Disk:} Tối thiểu 10GB cho Docker images, volumes và build artifacts.
\end{itemize}

\section*{A.2. Biên dịch backend}

Từ thư mục gốc của dự án, thực thi:

\begin{verbatim}
./mvnw clean package -DskipTests
\end{verbatim}

Lệnh này biên dịch toàn bộ 17 module Maven theo thứ tự phụ thuộc và đóng gói các service
thành file JAR thực thi. Khi cần chạy test Maven đầy đủ, có thể bỏ tham số
\texttt{-DskipTests} hoặc chạy theo từng module/suite để giảm tải tài nguyên.

\section*{A.3. Chạy backend và hạ tầng local}

\begin{verbatim}
docker compose up -d --build
\end{verbatim}

Lệnh trên khởi động PostgreSQL, Redis, Kafka, Elasticsearch, Keycloak, Mailpit, Prometheus,
Grafana, Zipkin, ba Spring Cloud service và 13 business service. Các lần chạy đầu tiên có
thể mất vài phút do cần build image và tải dependency.

\section*{A.4. Chạy frontend local}

Frontend nằm trong thư mục \path|frontend/|. Khi cần chạy storefront hoặc Admin Panel ở chế độ
development:

\begin{verbatim}
cd frontend
npm install
npm run dev
\end{verbatim}

Mặc định Next.js dùng \path|http://localhost:3000|. Nếu Grafana local cũng đang dùng cổng
3000, cần đổi cổng một trong hai thành phần hoặc dùng cấu hình production-like nơi Grafana
được remap sang host port \texttt{3001}.

\section*{A.5. Kiểm tra trạng thái local}

\begin{verbatim}
# Xem trạng thái container
docker compose ps

# Xem log Gateway hoặc order-service
docker compose logs -f api-gateway
docker compose logs -f order-service

# Kiểm tra health Gateway
curl http://localhost:8080/actuator/health/readiness

# Kiểm tra Eureka registry
curl -u eureka:eureka http://localhost:8761/eureka/apps
\end{verbatim}

\section*{A.6. Các URL local quan trọng}

\begin{table}[H]
\centering
\caption{Các URL truy cập hệ thống trên môi trường local}
\begin{compacttable}
\begin{tabularx}{\textwidth}{|L{3.1cm}|L{4.3cm}|Y|}
\hline
\textbf{Thành phần} & \textbf{URL} & \textbf{Ghi chú} \\ \hline
Frontend & \path|http://localhost:3000| & Storefront/Admin Panel khi chạy Next.js dev \\ \hline
API Gateway & \path|http://localhost:8080| & Cổng vào backend \\ \hline
Swagger UI & \path|http://localhost:8080/swagger-ui.html| & Tài liệu API qua Gateway \\ \hline
Eureka Dashboard & \path|http://localhost:8761| & user: eureka, pass: eureka \\ \hline
Keycloak Admin & \path|http://localhost:8180| & Quản trị realm \\ \hline
Grafana & \path|http://localhost:3000| hoặc \path|3001| & Phụ thuộc cấu hình chạy frontend/Grafana \\ \hline
Zipkin & \path|http://localhost:9411| & Distributed tracing \\ \hline
Prometheus & \path|http://localhost:9090| & Metrics \\ \hline
Mailpit & \path|http://localhost:8025| & Kiểm tra email \\ \hline
Elasticsearch & \path|http://localhost:9200| & Search engine \\ \hline
\end{tabularx}
\end{compacttable}
\end{table}

\section*{A.7. Chạy smoke test backend}

\begin{verbatim}
.test/scripts/smoke-backend.sh --suite maven
.test/scripts/smoke-backend.sh --suite infra
.test/scripts/smoke-backend.sh --suite auth-user
.test/scripts/smoke-backend.sh --suite catalog-search
.test/scripts/smoke-backend.sh --suite cart-voucher
.test/scripts/smoke-backend.sh --suite order-cod-review
.test/scripts/smoke-backend.sh --suite vnpay
.test/scripts/smoke-backend.sh --suite flash-sale
.test/scripts/smoke-backend.sh --suite security
\end{verbatim}

Kết quả final của các suite đã được lưu trong thư mục \texttt{.test/results} và được tổng hợp
ở Chương 5.

\section*{A.8. Chạy kiểm thử hiệu năng bằng k6}

Ví dụ chạy catalog soak hoặc checkout stress với gateway production-like:

\begin{lstlisting}[numbers=none,xleftmargin=0pt,breaklines=true]
k6 run -e BASE_URL=http://13.213.118.96:8080 .test/load/catalog-browse.js
k6 run -e BASE_URL=http://13.213.118.96:8080 .test/load/checkout-stress.js
\end{lstlisting}

Khi đưa số liệu vào báo cáo, chỉ dùng các artifact raw đã lưu trong \texttt{.test/results}
để tránh ghi số liệu không đối chiếu được.

\section*{A.9. Vận hành EC2 production-like}

Khi cần bật EC2 để demo:

\begin{lstlisting}[numbers=none,xleftmargin=0pt,breaklines=true]
source aws/config.env
aws ec2 start-instances --instance-ids $INSTANCE_ID
aws ec2 wait instance-running --instance-ids $INSTANCE_ID

# Doi 2-3 phut cho services khoi dong
ssh -i $SSH_KEY_PATH ubuntu@$ELASTIC_IP "docker compose ps"
\end{lstlisting}

Khi demo xong, nên tắt instance để tiết kiệm chi phí/credit:

\begin{lstlisting}[numbers=none,xleftmargin=0pt,breaklines=true]
source aws/config.env
aws ec2 stop-instances --instance-ids $INSTANCE_ID
\end{lstlisting}

\section*{A.10. Dừng hệ thống local}

\begin{verbatim}
# Dừng và giữ dữ liệu
docker compose down

# Dừng và xóa toàn bộ dữ liệu volume
docker compose down -v
\end{verbatim}

\end{document}
